\documentclass[a4paper]{article}
% Uncomment the following line to allow the usage of graphics (.png, .jpg)
\usepackage[pdftex]{graphicx}
% Allow the usage of utf8 characters
\usepackage[utf8]{inputenc}
\usepackage{enumerate}
\usepackage{amssymb}
\usepackage{colortbl}
\usepackage{icomma}
\usepackage[ngerman]{babel}
\usepackage{amsmath}
\usepackage{parskip}
\usepackage{booktabs}
\usepackage{siunitx}
\sisetup{locale=DE}
\usepackage{geometry}
\geometry{a4paper, top=15mm, left=15mm, right=15mm, bottom=15mm,
	headsep=10mm, footskip=12mm}
\usepackage{href-ul}
\hypersetup{
	colorlinks=true,
	linkcolor=blue,
	urlcolor=blue}
\begin{document}
	
	\title{Humorous Physics – The Global Warming Law}
	
	\section*{Formula of the Day}
	
	\[
	Q = c \cdot m \cdot \Delta \vartheta
	\]
	
	\begin{itemize}
		\item \( Q \): Heat energy in joules (J)
		\item \( c \): Specific heat capacity in \( \frac{\text{J}}{\text{kg} \cdot \text{K}} \)
		\item \( m \): Mass in kilograms (kg)
		\item \( \Delta \vartheta \): Temperature change in Kelvin (K) or degrees Celsius (°C)
	\end{itemize}
	
	\hrulefill
	
	\section*{1. The Lazy Cat and the Hot Water Bottle}
	
	Cat Leopold (mass: 6 kg) is lying on a hot water bottle containing 2 liters of water at room temperature (20 °C). The hot water bottle is heated to 60 °C so that his paws don't freeze.
	
	\begin{itemize}
		\item a) How much heat energy is needed to heat the water?
		\( c_\text{water} = 4{.}18\,\frac{\text{kJ}}{\text{kg}\cdot\text{K}} \)
		
		\vspace{1em}
		\( Q = \underline{\hspace{6cm}} \)
		
		\item b) Bonus: If the cat purrs and converts 20 \% of the energy into sound, how much remains as heat?
		
		\vspace{1em}
		\( Q_\text{heat} = \underline{\hspace{6cm}} \)
	\end{itemize}
	
	\hrulefill
	
	\section*{2. The Spaghetti Action }
	
	You want to boil 0{.}5 kg of water on the stove – from 15 °C to 100 °C. Your stove is in a bad mood and only delivers 80 \% of the energy used as actual heat.
	
	\begin{itemize}
		\item a) Calculate the energy required for ideal implementation.
		
		\vspace{1em}
		\( Q_\text{ideal} = \underline{\hspace{6cm}} \)
		
		\item b) Calculate how much energy you actually need to supply.
		
		\vspace{1em}
		\( Q_\text{real} = \underline{\hspace{6cm}} \)
	\end{itemize}
	
	\hrulefill
	
	\section*{3. The bathtub for astronauts}
	
	A space station simulates a cozy bath: 100 kg of water is heated from 18 °C to 38 °C. Energy is expensive there!
	
	\begin{itemize}
		\item a) Calculate the amount of heat required.
		
		\vspace{1em}
		\( Q = \underline{\hspace{6cm}} \)
		
		\item b) How many 1000-watt heating elements do you need if you want the bath to be warm in 10 minutes?
		
		(1 watt = 1 joule per second)
		
		\vspace{1em}
		\# heating elements = \underline{\hspace{6cm}}
	\end{itemize}
	
	\hrulefill
	
	\newpage
	
	\section*{4. The lukewarm declaration of love }
	
	A physics nerd in love brings his beloved a cup of tea with 250 ml of water. The tea should be heated to exactly 50 °C, but comes out of the tap at 15 °C. He carries the tea through the cold for 5 minutes, where 10% of the energy is lost.
	
	\begin{itemize}
		\item a) Calculate the amount of heat required without losses.
		
		\vspace{1em}
		\( Q_\text{ideal} = \underline{\hspace{6cm}} \)
		
		\item b) How much energy does he actually need to supply?
		
		\vspace{1em}
		\( Q_\text{real} = \underline{\hspace{6cm}} \)
	\end{itemize}
	
	\section*{5. The Hot Screw}
	
	A blacksmith heats an iron screw with a mass of 300 g from 20 °C to 180 °C.
	Specific heat capacity of iron: \( c = 450\, \frac{\text{J}}{\text{kg}\cdot\text{K}} \)
	
	\begin{itemize}
		\item a) Calculate the heat energy:
		
		\vspace{1em}
		\( Q = \underline{\hspace{6cm}} \)
		
		\item b) What would \( Q \) be if the screw weighed 600 g?
		
		\vspace{1em}
		\( Q = \underline{\hspace{6cm}} \)
	\end{itemize}
	
	\hrulefill
	
	\section*{6. The melting butter }
	
	A 250 g block of butter is heated from 6 °C to 40 °C.
	\( c_\text{Butter} = 2100\, \frac{\text{J}}{\text{kg}\cdot\text{K}} \)
	
	\begin{itemize}
		\item a) Calculate the absorbed heat energy:
		
		\vspace{1em}
		\( Q = \underline{\hspace{6cm}} \)
		
		\item b) What temperature would 125 g of butter reach if the same amount of energy were added?
		
		\vspace{1em}
		\( \Delta \vartheta = \underline{\hspace{6cm}} \)
	\end{itemize}
	
	\hrulefill
	
	\section*{7. The Fiery Brick }
	
	A brick (2 kg) is heated in the oven from 25 °C to 900 °C.
	\( c = 840\, \frac{\text{J}}{\text{kg}\cdot\text{K}} \)
	
	\begin{itemize}
		\item a) Calculate the absorbed heat energy:
		
		\vspace{1em}
		\( Q = \underline{\hspace{6cm}} \)
		
		\item b) Additional question: What role does color play in heating in the sun?
		
		\vspace{1em}
		\underline{\hspace{12cm}}
	\end{itemize}
	
	\hrulefill
	
\section*{8. The cheese cube in the toaster }

A cheese cube (50 g) is heated in the toaster from 10 °C to 60 °C.
\( c = 3100\, \frac{\text{J}}{\text{kg}\cdot\text{K}} \)

\begin{itemize}
	\item a) Calculate the required heat energy:
	
	\vspace{1em}
	\( Q = \underline{\hspace{6cm}} \)
	
	\item b) How long does a 150-watt toaster need for this?
	
	\vspace{1em}
	\( t = \underline{\hspace{6cm}} \)
\end{itemize}

\hrulefill

\section*{9. Sauna with wooden bench }

A pine block (1.5 kg) heats up from 25 °C to 75 °C.
\( c_\text{wood} = 1700\, \frac{\text{J}}{\text{kg}\cdot\text{K}} \)

\begin{itemize}
	\item a) Calculate the absorbed heat energy:
	
	\vspace{1em}
	\( Q = \underline{\hspace{6cm}} \)
	
	\item b) How much energy would an aluminum block of the same weight absorb?
	\( c_\text{Alu} = 900\, \frac{\text{J}}{\text{kg}\cdot\text{K}} \)
	
	\vspace{1em}
	\( Q = \underline{\hspace{6cm}} \)
\end{itemize}

\fbox{
	\begin{minipage}{0.5\textwidth}
		For the solution, please \href{https://www.okuyakl.de/phys/p9qcmL116/le116.pdf}{click here}or scan the QR code.\\
		
		More worksheets are available at
		
		\href{https://www.okuyakl.de}{www.okuyakl.de}
	\end{minipage}
	\hfill
	\begin{minipage}{0.4\textwidth} 
		\includegraphics[width=1.5 cm]{../../viecher/zwe03}
		\includegraphics[width=3 cm]{qre116}
		\includegraphics[width=2 cm]{../../viecher/afanticon1}
		
\end{minipage}}

\end{document}